\subsection{All declared runs and certificate costs}
The source-frozen execution retains all 225 timed records: 88 exact; 23 numerical; 2 state-limit; 29 time-limit; 83 tolerance. Status counts summarize record completeness, not a pooled success rate across unequal accuracies, budgets, or method definitions. The separate panels in Tables~\ref{tab:primary56}--\ref{tab:lattice56} retain their original targets. 200 runs produce rational certificates, including conservative fallback intervals from interrupted algorithms. 200 certificates have a successful independent check. 3 initial checks reach their twelve-second allowance; 3 unchanged certificates subsequently pass the separately timed extended check. No initial status is rewritten. The longest extended check takes 36.358 seconds. Numerical MIP returns remain separate from exact rational evidence.

The heterogeneous deficit method completes 27 of 30 requested primary and accuracy runs. 27 completed runs have an independently checked exact reference on the identical model. The maximum measured incumbent loss among those comparisons is 0.00111898. This observed loss is not the guaranteed upper-minus-lower interval. Interrupted resource computations and expensive independent checking remain part of the record; the faster convolution theorem does not make every complete proof pipeline inexpensive.

\subsection{Exact-target search and end-to-end screening}
The exact-target extension contains 18 price runs, of which 17 close at zero rational width. Table~\ref{tab:exact56} reports all nine two-second cases, with nodes, oracle calls, root gaps, fixed-book comparisons, depth and elapsed optimization time. The half-second counterparts remain in the raw records. Nine separate deterministic replays expose each tested price and required/forbidden set. They measure search mechanics, not replacement benchmark times. Across recorded price runs, maximum tree depth is 3 and the largest completed-oracle count is 45. Retained prices use at most 21 numerator bits and 20 denominator bits. The original R54 hard reduction instances, including failures, remain unchanged in the historical evidence. Neither root closure nor these observed bit lengths imply a polynomial worst-case search tree.

There are 18 paired screening comparisons under the same two-second total optimization allowance. 0 pairs improve total optimization-plus-serialization-plus-checking time; the median paired total-time ratio is 1.604. Table~\ref{tab:screen56} reports charges, deletions by eligibility category, final widths and both total times. Screening is charged for incumbent construction and forced-command bounds before the reduced solve. All unfavorable pairs remain visible. A smaller catalog or tighter bound is not itself an end-to-end speed improvement.

\subsection{Lattice scaling and implementation checks}
Among the eighteen zero-service lattice cases, 18 return an independently checked exact optimum. Cap-grid denominators are 8, 16 and 32; command budgets are two and four; all three seeds are retained. The table reports the actual resource-grid size, which also depends on probabilities and commanded-quantity granularity. The separate encoding stress uses denominators $2^8$, $2^{16}$ and $2^{32}$ and preserves any state-limit rejection rather than assigning it an optimum.

The exact regression covers 36 general models, 36 lattice models, 36 common-reward cap-count constructions, four strict $2q$ examples, and 200 convolutions with support holes. The joint cap--reward extension adds 36 genuinely heterogeneous reward models, each with two joint classes and independently checked original-policy reconstruction. Seven certificate mutation classes are rejected. A full twenty-thousand-bit rational certificate and a pre-allocation enormous-state rejection distinguish correct serialization from numerical tractability. These finite tests validate implementations; the universal assertions rely on the stated proofs. All operational data remain synthetic.
